\section{Simple comparison is not simple}
\begin{quotation}\textit{
Solving a quadratic at least looks as if it enough of a task to be worth
thinking about, even though at the start it does not appear to need a great
deal of thought. But there can be yet smaller activities where computer
arithetic has some subtlety. The one looked at here is just that of
comparing two numbers to decide which is bigger! By all rights that
should surely be utterly straightforward\ldots
}\end{quotation}

Suppose that in a computer program you have two variables. The first is
an integer \verb+i+ and here the case to be considered is when this
integer is 64-bits wide. The second variabe \verb+d+ is a double precison
float, and that too will fill 64-bits. In the past there were multiple
different ways in which floating point values ended up packed into
computer memory and much complication handling all the variations, but today
one can be assured that this floating point value is in a standard
form defined by the IEEE organization. The task is the trivial one of
deciding which of these two variables holds the larger numeric value.
One is tempted to write something as simple and direct as \verb+i > d+ as
the key test and expect everything to behave sensibly. But it might not!

Both computer integers and floating point values have limitations. So
first imagine trying to make the comparison by first turning both
values into integers. Well if it is large the floating point number can easily
have a value beyond the range of computer integers, and any attempt to
convert it will necessarily give trouble. In the languages C and C++ this
issue is expressed by the language standard declaring that an attempt to
convert an over-large floating value to an integer can lead to
``undefined behaviour'' and their expensive explanation of that allows
for absolutely any consequences, not at all limited to ones that lie
within the bounds of sanity. On the other hand if the floating point number's
value is more modest there may well be a fractional part and a conversion
to an integer necessarily loses this.

So perhaps the best scheme will be to start by converting the integer
to floating point? That can always be done and there are no fractions
anywhere to cause confusion. However beyond a certain limit the floating
point reporesentation can represent whole numbers but not all of them. To
be specific for the IEEE format under consideration here any integer
with magnitude up to $2^{53}$ can be converted to a double precision number
with no loss, but from there up to the biggest integers (close to $\pm 2^{63}$)
there has to be rounding. At the extreme one finds that adjacent double
precision numbers differ in value by 1024 (=$2^{10}$). Common programming
languages arrange that when an integer is mapped into floating point and
there needs to be rounding that works by selecting the nearest value.
So the double precision value you end up with may be either larger or
smaller than (or indeed sometimes equal to) the number you started with.
So you can not follow through to a simple reliable comparison!

Perhaps the reader might like to pause here and see if they can find
a neat way around all this.

One way to get started would be to handle a succession of straightforward
cases first. So if \verb+d+ is greater than $2^{63}$ or less than $-2^{63}$ it
is beyond the range of 64-bit integers and so depenings on its sign it
is unambiguously larger or smaller than whatever integer we have. It may
also be proper to check for the double precision value being a NaN (Not a
Number) because that would be unordered against any other numeric value.
If the floating point value is somewhat smaller then it might be converted
to an integer. If it has a fractional part that has numeric value less than
1.0 and so if the double value exceeds one greater than the integer or
is less then the integer minus one a result can be declared. Finally
one can obtain the fractional part of the floating point number by taking
the full value and subtracting what emerges from a a round trip that
converts it to an integer and then back. As always in such a detailed sequence
of steps i t will be necessary to think about edge cases. A particular one
is that if the floating point value is exactly equal to $2^{63}$ and attempt to
convert it to an integer would overflow. This write-up does not present a
full and clean version following this approach because in the following
paragraph presents an alterative but perhaps less obvious approach.

So for a second attempt the presentation here will be in the form of lines
of code, using C style annotations to show conversions from one type
to another. So a first step will be to make a value \verb+w+ by converting
the integer into a floating point value:
\begin{verbatim}
        double w = (double)i;
\end{verbatim}
Now we know that for large integers this value \verb+w+ may have to be either
larger than the accurate value of \verb+i+, but it will bn one of the two
possible floating point numbers one lower and one larger than the ideal.
The following picture suggests the scheme where large dots are values
in the floating point domain and small ones are those that 64-bit integers
can stand for. Positions on this line corresponding to one possibility
for \verb+i+ and \verb+w+ are marked. 
\begin{verbatim}
              ...O...............O...
                 ^    ^
                 w    i
\end{verbatim}
Any value of \verb+d+ that is strictly greater than w must be strictly greater
than \verb+i+ because of the spacing between values representable in
floating point. Similarly anything strictly less than \verb+w+ is of
necessity also less than \verb+i+. That means that a continuation of the code
we need might be
\begin{verbatim}
        if d > w then say "greater"
        if d < w then say "less"
\end{verbatim}
and now the only case to worry about will be when \verb+d+ and \verb+w+
happen to be equal (which may be a tolerably uncommon case).
Before continuing it may be ueful to consider the case if the only test
that had to be made is whether \verb+d+ is strictly greater than \verb+i+.
In that situation if \verb+d+ is a NaN it is proper to report a ``no it
is not greater''. and this is turns out to be neatly incorporated by
adjusting the second of the above lines to be
\begin{verbatim}
        if d != w then say "less or unordered"
\end{verbatim}

To continue with the case where \verb+d+ is equal to \verb+w+. A first
thought is that \verb+w+ came from an integer and so is not badly outside the
range of integers. Well actually if it has the value $9223372036854775808.0$
which is exactly $2^{63}$ and can have arisen from converting and hence
rounding up integer that started off a little less than that. This is the
{\em only} special case! so it gets handled using 
\begin{verbatim}
        if w = 9223372036854775808.0 then say "greater"
\end{verbatim}
Because \verb+w+ and \verb+d+ are now known to be equal it would be
valid to check either of them. Finally it is possible to end up with
code that is much like one of the very earliest thoughts, because
now we are assured that the conversion from double to integer will be
in range and will be exact:
\begin{verbatim}
        return what you get from the comparison (integer)d > i
\end{verbatim}

About the only reasonable response to this is amazement first that such
a trivial looking task ends up so messy to cope with, but then that there
is a solution that is concise and liable to be reasonably efficient!

